\section{Related Work}
\textbf{Feature extraction with pretrained models.}
Using pretrained neural network models for extracting visual features \old{is an approach that} has stood the test of time \old{ever} since the AlexNet \citep{alexnet} CNN model pretrained on ImageNet-1k~\citep{russakovsky2015imagenet}. 
More recent models have upgraded the same setup with modern architectures, such as ResNets (used in, \emph{e.g.}, DETR,~\citealp{carion2020end}) or even Vision Transformers. 
As Transformers are easily able to handle different modalities during training, off-the-shelf backbones are now commonly trained on label supervision (\emph{e.g.}, DeiT-III on ImageNet-22k,~\citealp{touvron2022deit}) or text supervision (e.g., CLIP~\citep{radford2021learning}), providing strong \textit{visual foundation models}, scaling well with model sizes, and enabling excellent performance on a variety of tasks including detection \citep{carion2020end} and segmentation \citep{zheng2021rethinking,kirillov2023segment}.
In this context, supervision relies on annotations in the form of labels or text alignment; the dataset biases \citep{torralba2014dataset} are not well characterized, yet they drive learning and shape the learned models. 
An alternative approach consists of not using supervision and letting the models learn from the data via a pretext task that is designed to require understanding the content of images~\citep{doersch2015unsupervised}. 
This self-supervised learning paradigm was explored in multiple methods using Vision Transformers: MAE \citep{he2021masked} trains a model at reconstructing pixel values of hidden areas of an image and then applies fine-tuning to address a new task. 
With a different approach, the self-distillation family of methods \citep{he2020momentum, caron2021emerging, zhou2021ibot} showcase strong performance using frozen backbones, allowing for more robustness to domain shifts for task-specific downstream models.
In this work, we focused the analysis on self-supervised learning, and more specifically on the DINOv2 approach \citep{oquab2023dinov2}, which has shown to be particularly effective for learning local features.
We showed that despite excellent benchmark scores, DINOv2 features exhibit undesirable artifacts and that correcting these artifacts in the learning process allows for further improvements in the benchmark performances.
These phenomenon is even more surprising as DINOv2 builds upon DINO~\citep{caron2021emerging}, which does not show signs of artifacts.
We then further showed that the correction techniques hold for supervised paradigms by testing on DeiT-III and OpenCLIP.

\textbf{Additional tokens in transformers.}
Extending the transformer sequence with special tokens was popularized in BERT~\citep{devlin2018bert}. 
However, most approaches add new tokens either to provide the network with new information as for example \texttt{[SEP]} tokens in BERT, provide opportunity to spend more computation on the input as seen with the tape tokens in AdaTape~\citep{xue2023adaptive}, or to gather information in these tokens, and use their output value as an output of the model:
for classification, as \texttt{[CLS]} tokens in BERT and ViT \citep{dosovitskiy2020image};
for generative learning, as \texttt{[MASK]} in BERT and BEiT~\citep{bao2021beit};
for detection, as object queries in DETR~\citep{carion2020end}, detection tokens in YOLOS~\citep{fang2021you}, and ViDT~\citep{song2021vidt};
or for accumulating information from possibly multiple modalities before decoding, as latent token arrays in Perceivers ~\citep{jaegle2021perceiver,Jaegle2021PerceiverIO}.
Different to these works, the tokens we add to the sequence add no information, and their output value is not used for any purpose. 
They are simply registers where the model can learn to store and retrieve information during the forward pass.
The Memory Transformer~\citep{burtsev2020memory}, closer to our work, presents a simple approach to improve transformer models using memory tokens added to the token sequence, improving translation performance. 
In follow-up work, \cite{bulatov2022recurrent} address complex copy-repeat-reverse tasks. 
\cite{sandler2022fine} extend this line to the vision domain for fine-tuning but observe that such tokens do not transfer well across tasks.
In contrast, we do not perform fine-tuning and employ additional tokens during pretraining to improve the features obtained for all tasks downstream.
More importantly, our study contributes the following new insight in Sec.~\ref{sec:problem}: the mechanism implemented through memory tokens already appears naturally in Vision Transformers; our study shows that \textit{such tokens allow us not to create but to isolate this existing behavior}, and thus avoid collateral side-effects.

\textbf{Attention maps of vision transformers.}
Visualising the attention map from \texttt{[CLS]} token to patch tokens was popularized in DINO~\citep{caron2021emerging}.
It was shown there that the attention maps of DINO were clean of artifacts, as opposed to the attention maps of previous vision transformers.
Other works have since reported interesting attention maps using various techniques: by modifying the optimisation procedure~\citep{chen2021vision}, by steering the attention scores towards useful image parts~\citep{shi2023toast}, by modifying the architecture of the transformer layers~\citep{yu2023emergence}, or by introducing a learnable pooling to produce the \texttt{[CLS]} token~\citep{psomas2023keep}.